\documentclass{article}
\usepackage{graphicx} % Required for inserting images

% Added stuff
\usepackage{amsmath}
\usepackage{amssymb}
\newcommand{\N}{\mathbb{N}}
\newcommand{\Z}{\mathbb{Z}}
\newcommand{\Q}{\mathbb{Q}}
\newcommand{\R}{\mathbb{R}}
\newcommand{\C}{\mathbb{C}}
\newcommand{\bigp}[1]{\left(#1\right)}
\newcommand{\bigb}[1]{\left[#1\right]}
\newcommand{\bigabs}[1]{\left|#1\right|}
\newcommand{\summation}[1][0]{\sum_{n=#1}^\infty}
\newcommand{\summationk}[1][0]{\sum_{k=#1}^\infty}
\newcommand{\summationm}[1][0]{\sum_{m=#1}^\infty}
\newcommand{\eps}{\varepsilon}
\renewcommand{\emptyset}{\varnothing}

\title{p-adic Analysis 2}
\author{Siddharth K}
\date{April 2026}

\begin{document}

\maketitle

\section*{Problem 1}
Repeatedly solving $7a_n \equiv 1 \pmod{3^n}$ gives us
$$\frac{1}{7} = \ldots 01212010212011_3.$$

\section*{Problem 2}
Yes, since $5 \not\equiv 0 \pmod{11}$, there exists $5^{-1}$. Letting $x = 5^{-1} \cdot 3$, we have that $5x = 3$.

\section*{Problem 3}
We need $x^2 \equiv -1 \pmod{5}$, which is satisfied by any $x \equiv 2, 3 \pmod{5}$. Now, assume for some $n \ge 1$ we found $x_n$ where
$$x_n^2 \equiv -1 \pmod{n}.$$
We want $x_{n + 1} = x_n + a \cdot 5^n$ for some $a \in \{0, 1, 2, 3, 4\}$. So
$$x_{n + 1}^2 = x_n^2 + 2x_na5^n + a^25^{2n}.$$
Taking this $\bmod 5^{n + 1}$ gives us
$$x_n^2 + 2x_na5^n \equiv -1 \pmod{5^{n + 1}}.$$
Since $x_n^2 \equiv -1 \pmod{5^n}$, we can write $x_n^2 + 1 = c\cdot 5^n$ for some $c \in \Z$. Then we get
$$c\cdot 5^n + 2x_na5^n \equiv 0 \pmod{5^{n + 1}},$$
$$c + 2x_na \equiv 0 \pmod{5},$$
which has a unique solution $a \pmod{5}$.

The first 10 digits are $2121342303$.

There is no solution in $\Z_7$ because all the squares $\bmod 7$ are
$$0^2 \equiv 0, 1^2 \equiv 1, 2^2 \equiv 4, 3^2 \equiv 2, 4^2 \equiv 2, 5^2 \equiv 4, 6^2 \equiv 1,$$
none of which are $6$.


\section*{Problem 11}
These are just the rational numbers that have nonnegative $p$-adic valuation, which are rationals of the form $a/b$ where $p \nmid b$.

\section*{Problem 12}
Let $x \in \Q_p^\times$. There is a unique $n \in \Z$ such that $x = p^ny$ with $y \in \Z_p^\times$. So the problem reduces to writing $y = ru$ with $r \in \{1, \ldots, p - 1\}$ and $u \in 1 + p\Z_p$.

Since $y \in \Z_p^\times$, there is a unique $r \in \{1, \ldots, p - 1\}$ with
$$y \equiv r \pmod{p}.$$
Then $r$ is a p-adic unit as well, so $r^{-1} \in \Z_p$. Define
$$u = r^{-1}y,$$
so that $y = ru$ and $x = p^n ru$. Since $y \equiv r \pmod{p}$, we get
$$u \equiv 1 \pmod{p},$$
as desired. This is unique because both $n$ and $r$ were forced to be unique.

\section*{Problem 15}
The set $\Q$ is already complete with respect to the trivial absolute value because Cauchy sequences are eventually constant rational numbers, so they always converge to rational numbers.

\section*{Problem 16}
We have that
$$\prod_p |x|_p = \prod_p p^{-v_p(x)} = \bigp{\prod_p p^{v_p(x)}}^{-1}.$$
But by definition,
$$\prod_p p^{v_p(x)} = |x|$$
since it is the exact prime factorization of $|x|$. Since $|x| = |x|_\infty$,
$$|x|_\infty \prod_p |x|_p = |x|_\infty \cdot |x|^{-1} = 1.$$

\section*{Problem 17}
Assume $x = r^2$ for some $r \in \Q$. Since $r \in \R$ and $r \in \Q_p$, $x = r^2$ is a perfect square in $\R$ and $\Q_p$.

Now, assume $x$ is a square in $\Q_p$ for every $p$, and also a square in $\R$ (which forces it to be positive and finite).
Fix a prime $p$. Since $x$ is a square in $\Q_p$, there exists $y_p \in \Q_p$ such that $x = y_p^2$, and therefore $v_p(x) = 2v_p(y)$. So $v_p(x)$ is even for all $p$. Now, write
$$x = \prod_p p^{e_p}.$$
Since each $e_p$ is even, we can let $f_p = e_p/2$ and write
$$x = \bigp{\prod_p p^{f_p}}^2,$$
so $x$ is a perfect square in $\Q$.

\section*{Problem 18}
\subsection*{Part (b)}
Consider
$$\R = \bigcup_{n = 1}^\infty (-n, n).$$
Any finite subcover is of the form
$$(-n_1, n_1)\cup \ldots\cup(-n_k, n_k) = (-N, N)$$
where $N = \max_i(n_i)$, and this doesn't cover $\R$.

\subsection*{Part (c)}
First assume $A$ is compact. It has to be unbounded because otherwise,
$$U_n = (-n, n), \quad n \ge 1$$
would form an open cover of $A$ with no finite subcover. To show $A$ is closed, let $x \in \R \setminus A$. For each $a \in A$, set
$$U_a = \left(a - \frac{1}{2}|x - a|,\; a + \frac{1}{2}|x - a|\right).$$
Then $\{X_a\}_{a \in A}$. By compactness, only finitely many are required:
$$A \subseteq U_{a_1} \cup \cdots \cup U_{a_n}.$$
Let
$$\eps = \min_i \frac{|x - a_i|}{2} > 0.$$
Then $(x - \eps, x + e)$ doesn't intersect with any of the $U_{a_i}$, so it misses $A$. So $R \setminus A$ is open, and therefore $A$ is closed.

\end{document}